% =====================================================================
% Part 1 -- Learning-theory cluster for the shift-invariance / robustness primer.
% Three sections + closing bridge box:
%   1. Margins and the max-margin classifier (SVM), convex-hull view.
%   2. What gradient descent actually learns -- implicit bias.
%   3. NTK and the lazy-vs-rich dichotomy.
% This file is meant to be \input into a host document that already loads
% amsmath/amsthm and defines the theorem environments
%   theorem, lemma, corollary, proposition, definition, example, remark
% (exactly as in primer.tex). No \documentclass / \begin{document} here.
% Every concept is attributed with \citep{...}; the key list is at the end.
% =====================================================================

\section{Margins and the max-margin classifier}
\label{sec:margins}

Everything in this primer eventually comes back to a single geometric
quantity: \emph{how far is a data point from the surface where the
classifier changes its mind?} That distance is called the \emph{margin},
and a classifier that makes it as large as possible is a
\emph{max-margin classifier}, the simplest of which is the
(hard-margin) \emph{support vector machine}, or SVM
\citep{boser1992, cortes1995, vapnik}. We build the idea up slowly,
prove the one formula that every textbook quietly assumes, and then
look at the same object from a second angle (closest points of two
convex hulls) that will make the later robustness story intuitive.

\subsection{Hyperplanes and two notions of margin}

We work in $\mathbb{R}^d$. Our data are pairs $(\mathbf{x}_i, y_i)$ with
inputs $\mathbf{x}_i \in \mathbb{R}^d$ and binary labels
$y_i \in \{+1, -1\}$. A \emph{linear classifier} predicts the label of a
new point $\mathbf{x}$ by looking at the sign of an affine function
$\mathbf{w}\cdot\mathbf{x} + b$, where $\mathbf{w}\in\mathbb{R}^d$ is a
\emph{weight vector} and $b\in\mathbb{R}$ is a \emph{bias} (also called an
\emph{intercept} or \emph{threshold}). Here $\mathbf{w}\cdot\mathbf{x}
= \sum_{j=1}^d w_j x_j$ is the ordinary dot product. The prediction is
$+1$ when $\mathbf{w}\cdot\mathbf{x}+b>0$ and $-1$ when it is negative.

\begin{definition}[Hyperplane]
A \emph{hyperplane} in $\mathbb{R}^d$ is the set of points
$\{\mathbf{x} : \mathbf{w}\cdot\mathbf{x} + b = 0\}$ for some nonzero
$\mathbf{w}$ and scalar $b$. It is a flat set of dimension $d-1$: a line
in $\mathbb{R}^2$, an ordinary plane in $\mathbb{R}^3$. The vector
$\mathbf{w}$ is \emph{normal} (perpendicular) to the hyperplane, and the
hyperplane is the \emph{decision boundary} of the classifier: the place
where the score $\mathbf{w}\cdot\mathbf{x}+b$ is exactly zero, so the
predicted label flips.
\end{definition}

Why is $\mathbf{w}$ perpendicular to the hyperplane? Take any two points
$\mathbf{x}$ and $\mathbf{x}'$ that lie on it. Both satisfy the equation,
so $\mathbf{w}\cdot\mathbf{x}+b = 0$ and $\mathbf{w}\cdot\mathbf{x}'+b = 0$.
Subtracting, $\mathbf{w}\cdot(\mathbf{x}-\mathbf{x}') = 0$. The vector
$\mathbf{x}-\mathbf{x}'$ points along the hyperplane (it is the difference
of two of its points), and it is orthogonal to $\mathbf{w}$; since
$\mathbf{x},\mathbf{x}'$ were arbitrary, $\mathbf{w}$ is orthogonal to
every direction inside the hyperplane.

There are two different things one might call ``the margin,'' and keeping
them apart is the single most common source of confusion, so we name them
both \citep{burges1998}.

\begin{definition}[Functional and geometric margin]
For a labelled point $(\mathbf{x}_i, y_i)$ and a classifier
$(\mathbf{w}, b)$, the \emph{functional margin} is
\[
  \hat{\gamma}_i \;=\; y_i\,(\mathbf{w}\cdot\mathbf{x}_i + b).
\]
The \emph{geometric margin} is the functional margin divided by the length
of the weight vector,
\[
  \gamma_i \;=\; \frac{y_i\,(\mathbf{w}\cdot\mathbf{x}_i + b)}{\|\mathbf{w}\|},
\]
where $\|\mathbf{w}\| = \sqrt{\mathbf{w}\cdot\mathbf{w}}$ is the Euclidean
norm.
\end{definition}

The functional margin is positive exactly when the point is classified
correctly (the label $y_i$ and the score share a sign). But it is
\emph{not} a geometric distance: if we multiply both $\mathbf{w}$ and $b$
by $10$, the decision boundary does not move at all (the equation
$10\mathbf{w}\cdot\mathbf{x}+10b=0$ has the same solution set), yet the
functional margin grows tenfold. The geometric margin fixes this: it is
\emph{scale-invariant} (replacing $(\mathbf{w},b)$ by $(c\mathbf{w},cb)$
for $c>0$ leaves it unchanged because the $c$ in the numerator cancels the
$c$ in $\|c\mathbf{w}\|=c\|\mathbf{w}\|$), and as we now prove it equals the
honest perpendicular distance from $\mathbf{x}_i$ to the boundary.

\subsection{The point-to-hyperplane distance, with proof}

Burges' classic tutorial states the distance from a point to a hyperplane
as a fact and moves on \citep{burges1998}. Because this formula is the
backbone of margins, certificates and attacks alike, we supply the proof
in full.

\begin{theorem}[Point-to-hyperplane distance]
\label{thm:dist}
Let $H = \{\mathbf{x} : \mathbf{w}\cdot\mathbf{x} + b = 0\}$ with
$\mathbf{w}\neq \mathbf{0}$, and let $\mathbf{x}_0\in\mathbb{R}^d$ be any
point. The Euclidean distance from $\mathbf{x}_0$ to $H$ is
\[
  \operatorname{dist}(\mathbf{x}_0, H)
  \;=\; \frac{|\mathbf{w}\cdot\mathbf{x}_0 + b|}{\|\mathbf{w}\|}.
\]
\end{theorem}

\begin{proof}
The distance from $\mathbf{x}_0$ to the set $H$ is by definition the
smallest value of $\|\mathbf{x}_0 - \mathbf{x}\|$ over all
$\mathbf{x}\in H$. We will (i) guess the closest point by walking from
$\mathbf{x}_0$ straight toward $H$ along the normal direction, and
(ii) prove no point of $H$ is nearer.

\emph{Step 1 (a unit normal).}
Set $\mathbf{n} = \mathbf{w}/\|\mathbf{w}\|$, the normal $\mathbf{w}$
rescaled to length one, so $\mathbf{n}\cdot\mathbf{n} = 1$.

\emph{Step 2 (walk along the normal).}
Consider the point obtained by moving a signed amount $t\in\mathbb{R}$ from
$\mathbf{x}_0$ along $\mathbf{n}$,
\[
  \mathbf{p}(t) \;=\; \mathbf{x}_0 - t\,\mathbf{n}.
\]
We want the $t$ that lands $\mathbf{p}(t)$ on $H$, i.e.\ that makes
$\mathbf{w}\cdot\mathbf{p}(t)+b = 0$.

\emph{Step 3 (solve for $t$).}
Substitute and expand, using linearity of the dot product in the variable
$\mathbf{p}$,
\begin{align*}
  \mathbf{w}\cdot\mathbf{p}(t) + b
  &= \mathbf{w}\cdot(\mathbf{x}_0 - t\,\mathbf{n}) + b \\
  &= \mathbf{w}\cdot\mathbf{x}_0 \;-\; t\,(\mathbf{w}\cdot\mathbf{n}) \;+\; b.
\end{align*}
Now $\mathbf{w}\cdot\mathbf{n} = \mathbf{w}\cdot(\mathbf{w}/\|\mathbf{w}\|)
= \|\mathbf{w}\|^2/\|\mathbf{w}\| = \|\mathbf{w}\|$. Setting the expression
to zero and solving for $t$,
\[
  \mathbf{w}\cdot\mathbf{x}_0 - t\,\|\mathbf{w}\| + b = 0
  \quad\Longrightarrow\quad
  t^\star \;=\; \frac{\mathbf{w}\cdot\mathbf{x}_0 + b}{\|\mathbf{w}\|}.
\]
So the foot of the perpendicular is $\mathbf{p}(t^\star)\in H$, and the
length of the step we took is
\[
  \|\mathbf{x}_0 - \mathbf{p}(t^\star)\| = \|t^\star\,\mathbf{n}\|
  = |t^\star|\,\|\mathbf{n}\| = |t^\star|
  = \frac{|\mathbf{w}\cdot\mathbf{x}_0 + b|}{\|\mathbf{w}\|}.
\]

\emph{Step 4 (no point of $H$ is nearer).}
Let $\mathbf{x}\in H$ be arbitrary, so $\mathbf{w}\cdot\mathbf{x}+b = 0$.
Write the gap from $\mathbf{x}_0$ as
$\mathbf{x}_0 - \mathbf{x} = \big(\mathbf{x}_0-\mathbf{p}(t^\star)\big)
+ \big(\mathbf{p}(t^\star)-\mathbf{x}\big)$. The first piece equals
$t^\star\mathbf{n}$ (a multiple of the normal). The second piece,
$\mathbf{p}(t^\star)-\mathbf{x}$, is the difference of two points of $H$,
hence (as shown above) orthogonal to $\mathbf{w}$ and so to $\mathbf{n}$.
The two pieces are therefore orthogonal, and Pythagoras gives
\[
  \|\mathbf{x}_0 - \mathbf{x}\|^2
  = \|t^\star\mathbf{n}\|^2 + \|\mathbf{p}(t^\star)-\mathbf{x}\|^2
  \;\ge\; \|t^\star\mathbf{n}\|^2 = |t^\star|^2,
\]
with equality iff $\mathbf{p}(t^\star)=\mathbf{x}$. Hence the minimum
distance is exactly $|t^\star| = |\mathbf{w}\cdot\mathbf{x}_0+b|/\|\mathbf{w}\|$.
\end{proof}

Two corollaries we will reuse constantly. First, the distance from the
hyperplane to the \emph{origin} (set $\mathbf{x}_0 = \mathbf{0}$) is
$|b|/\|\mathbf{w}\|$. Second, comparing with the geometric margin:
$\gamma_i = y_i(\mathbf{w}\cdot\mathbf{x}_i+b)/\|\mathbf{w}\|$ is precisely
the signed distance of $\mathbf{x}_i$ to the boundary, positive when the
point is on its correct side. \emph{The geometric margin is a distance.}

\subsection{Canonical scaling, support vectors, and the hard-margin SVM}

Because the geometric margin does not change under rescaling, we are free
to fix the scale of $(\mathbf{w}, b)$ however is most convenient. The
convention everyone uses is the \emph{canonical} one
\citep{burges1998}: rescale so that the closest correctly classified
points sit at functional margin exactly one. Concretely we require

\begin{equation}
\label{eq:canonical}
  y_i\,(\mathbf{w}\cdot\mathbf{x}_i + b) \;\ge\; 1 \qquad\text{for all }i,
\end{equation}
with equality achieved by the nearest points. Under this convention the
two ``rails'' $\mathbf{w}\cdot\mathbf{x}+b = +1$ and
$\mathbf{w}\cdot\mathbf{x}+b = -1$ are parallel to the boundary and sit
on either side of it. By Theorem~\ref{thm:dist} each rail is at distance
$1/\|\mathbf{w}\|$ from the boundary, so the total width of the
separating slab is

\[
  \text{margin} \;=\; d_+ + d_- \;=\; \frac{1}{\|\mathbf{w}\|} + \frac{1}{\|\mathbf{w}\|}
  \;=\; \frac{2}{\|\mathbf{w}\|},
\]

where $d_+$ ($d_-$) is the distance to the nearest positive (negative)
point \citep{burges1998}.

\begin{definition}[Support vectors]
The \emph{support vectors} are the training points that actually touch a
rail, i.e.\ for which equality holds in \eqref{eq:canonical}:
$y_i(\mathbf{w}\cdot\mathbf{x}_i+b)=1$. They are the points lying exactly
on the slab boundary; the solution is determined entirely by them, and
moving or deleting a non-support point (one strictly inside its own half)
does not change the classifier.
\end{definition}

We want the widest slab, i.e.\ the largest $2/\|\mathbf{w}\|$. Since
$2/\|\mathbf{w}\|$ is a strictly decreasing function of $\|\mathbf{w}\|$,
maximizing the margin is the same as minimizing $\|\mathbf{w}\|$, and
equivalently $\tfrac12\|\mathbf{w}\|^2$ (a smooth, convex stand-in). This
is the \emph{hard-margin SVM} \citep{boser1992, vapnik, burges1998}:

\begin{equation}
\label{eq:svm-primal}
  \min_{\mathbf{w}, b}\ \tfrac{1}{2}\|\mathbf{w}\|^2
  \qquad\text{subject to}\qquad
  y_i\,(\mathbf{w}\cdot\mathbf{x}_i + b) \ge 1 \ \ \forall i.
\end{equation}

\begin{proposition}[Max-margin $\Longleftrightarrow$ minimum norm]
\label{prop:maxmargin}
Among all $(\mathbf{w}, b)$ satisfying the canonical constraints
\eqref{eq:canonical}, the one with the largest geometric margin is exactly
the minimizer of $\tfrac12\|\mathbf{w}\|^2$ in \eqref{eq:svm-primal}.
\end{proposition}
\begin{proof}
The margin under canonical scaling is $2/\|\mathbf{w}\|$. The map
$r \mapsto 2/r$ is strictly decreasing on $r>0$, so its maximum over the
feasible set is attained precisely where $\|\mathbf{w}\|$ (hence
$\tfrac12\|\mathbf{w}\|^2$) is smallest. The two optimizations have the
same minimizer.
\end{proof}

Problem \eqref{eq:svm-primal} is a convex quadratic program with linear
constraints, so it has a unique global solution $\mathbf{w}$; the
\emph{Karush--Kuhn--Tucker} (KKT) conditions characterize it. We pause to
define KKT cleanly because it returns in Sections~\ref{sec:implicit} and is
the language all of these results speak.

\begin{definition}[KKT point]
\label{def:kkt}
Consider $\min_{\mathbf{z}} f(\mathbf{z})$ subject to inequality
constraints $g_i(\mathbf{z})\ge 0$. A point $\mathbf{z}^\star$ is a
\emph{KKT point} if there exist multipliers $\alpha_i \ge 0$ (one per
constraint) such that
\begin{enumerate}[label=(\roman*),nosep]
\item \emph{stationarity:}
  $\nabla f(\mathbf{z}^\star) = \sum_i \alpha_i \nabla g_i(\mathbf{z}^\star)$;
\item \emph{primal feasibility:} $g_i(\mathbf{z}^\star)\ge 0$ for all $i$;
\item \emph{dual feasibility:} $\alpha_i \ge 0$ for all $i$;
\item \emph{complementary slackness:} $\alpha_i\, g_i(\mathbf{z}^\star) = 0$
  for all $i$.
\end{enumerate}
For a convex problem the KKT conditions are necessary \emph{and}
sufficient for a global minimum; for a non-convex problem they are only
necessary (a KKT point need not be a global minimizer). Intuitively, a
multiplier $\alpha_i$ is the ``price'' of constraint $i$, and
complementary slackness says we only pay for constraints that are tight
($g_i = 0$).
\end{definition}

Applying Definition~\ref{def:kkt} to the SVM \eqref{eq:svm-primal} with
$g_i(\mathbf{w},b) = y_i(\mathbf{w}\cdot\mathbf{x}_i+b)-1$, complementary
slackness reads
\[
  \alpha_i\,\big[\,y_i(\mathbf{w}\cdot\mathbf{x}_i+b)-1\,\big] = 0.
\]
So $\alpha_i>0$ forces $y_i(\mathbf{w}\cdot\mathbf{x}_i+b)=1$: the points
with nonzero multiplier are exactly the support vectors. Stationarity in
$\mathbf{w}$ gives $\mathbf{w} = \sum_i \alpha_i y_i \mathbf{x}_i$, so the
optimal weight vector is a nonnegative combination of the support vectors
alone. \emph{Remember this: the SVM solution is built only from the points
that sit on the margin.} Section~\ref{sec:implicit} will show plain
gradient descent rediscovers exactly this same combination on its own.

\subsection{The convex-hull view: closest points of two hulls}

There is a second picture, due to Bennett and Bredensteiner
\citep{bennett2000}, that many people find more intuitive than the
algebra and which previews the robustness story. It needs one definition.

\begin{definition}[Convex hull]
The \emph{convex hull} $\operatorname{conv}(A)$ of a finite set
$A = \{\mathbf{a}_1, \dots, \mathbf{a}_m\}$ is the set of all
\emph{convex combinations} of its points,
\[
  \operatorname{conv}(A) = \Big\{\ \textstyle\sum_{i} u_i \mathbf{a}_i \ :\
  u_i \ge 0,\ \sum_i u_i = 1 \ \Big\}.
\]
It is the smallest convex set containing $A$: in the plane, the rubber
band stretched around the points.
\end{definition}

Suppose the two classes are linearly separable. Let $A$ be the positive
points and $B$ the negative points. Bennett's theorem says the max-margin
hyperplane is determined by the two closest points of the two hulls.

\begin{theorem}[Closest points of the hulls \citep{bennett2000}]
\label{thm:hulls}
For separable classes, let
\[
  (\mathbf{c}^\star, \mathbf{d}^\star) = \arg\min_{\mathbf{c}\in\operatorname{conv}(A),\,
    \mathbf{d}\in\operatorname{conv}(B)} \ \tfrac{1}{2}\|\mathbf{c}-\mathbf{d}\|^2
\]
be the closest pair of points, one in each hull. Then the max-margin
hyperplane is the perpendicular bisector of the segment
$[\mathbf{c}^\star, \mathbf{d}^\star]$, its normal is
$\mathbf{w}\propto \mathbf{c}^\star - \mathbf{d}^\star$, and the margin
(slab width) equals the distance between the hulls,
$\|\mathbf{c}^\star - \mathbf{d}^\star\|$.
\end{theorem}

\begin{proof}[Proof skeleton]
The argument is by contradiction on optimality
\citep{bennett2000}. (a) \emph{The normal must be along
$\mathbf{c}^\star-\mathbf{d}^\star$.} If the optimal separating slab's
normal were \emph{not} parallel to the segment joining the closest hull
points, then the two parallel supporting planes (one touching each hull)
would not be as far apart as possible: tilting the slab to align with the
segment would increase the gap, contradicting maximality. (b) \emph{The
segment is orthogonal to the supporting planes.} If it were not, then the
foot of one endpoint could be slid along its hull to a strictly nearer
point, contradicting that $\mathbf{c}^\star,\mathbf{d}^\star$ are the
closest pair. Conditions (a) and (b) together force the boundary to be the
perpendicular bisector of $[\mathbf{c}^\star,\mathbf{d}^\star]$ with normal
$\mathbf{c}^\star-\mathbf{d}^\star$; the slab width is then the length of
that segment. The full statement is that the ``closest points of the
hulls'' problem is the Wolfe dual of the ``push two parallel supporting
planes apart'' problem, so they have the same solution
\citep{bennett2000}.
\end{proof}

The geometric reading is worth keeping: \emph{maximizing the margin is the
same as finding the two nearest points of the two clouds and bisecting the
segment between them.} (When classes overlap, Bennett shrinks each hull by
capping every weight $u_i\le 1/K$ -- the \emph{reduced convex hull} -- and
shows this softening of the hulls is the dual of allowing margin
violations, the soft-margin SVM of \citet{cortes1995}. We will not need
the soft-margin machinery, but it is good to know the picture extends.)

\subsection{Two worked examples and an exercise}

\begin{example}[Two points, margin $2$]
\label{ex:twopoint}
Take a single positive point $\mathbf{x}_+ = (2,0)$ with $y=+1$ and a
single negative point $\mathbf{x}_- = (0,0)$ with $y=-1$ in $\mathbb{R}^2$.
By symmetry the boundary is the vertical line halfway between them,
$x_1 = 1$. Let us verify both ways.

\emph{Canonical algebra.} The boundary is normal to the $x_1$-axis, so
$\mathbf{w}=(w_1,0)$. The canonical constraints with equality at both
points read $w_1\cdot 2 + b = +1$ and $-(w_1\cdot 0 + b) = +1$, i.e.\
$b=-1$ and $2w_1 - 1 = 1$, giving $w_1 = 1$. Then $\|\mathbf{w}\| = 1$ and
the margin is $2/\|\mathbf{w}\| = 2$. The boundary
$\mathbf{w}\cdot\mathbf{x}+b = x_1 - 1 = 0$ is indeed $x_1=1$.

\emph{Convex-hull view.} Each ``hull'' is a single point:
$\mathbf{c}^\star=(2,0)$, $\mathbf{d}^\star=(0,0)$. The normal is
$\mathbf{w}\propto\mathbf{c}^\star-\mathbf{d}^\star = (2,0)$, the bisector
is $x_1=1$, and the margin is
$\|\mathbf{c}^\star-\mathbf{d}^\star\| = 2$. The two pictures agree.
\end{example}

\begin{example}[Four facing segments, margin $3$]
\label{ex:foursegment}
Let $A = \{(3,1),(3,-1)\}$ ($y=+1$) and $B = \{(0,1),(0,-1)\}$ ($y=-1$).
Each hull is now a vertical segment: $\operatorname{conv}(A)$ is the
segment from $(3,-1)$ to $(3,1)$, and $\operatorname{conv}(B)$ the segment
from $(0,-1)$ to $(0,1)$. The two segments face each other; the closest
distance between them is the horizontal gap, $\|(3,y)-(0,y)\| = 3$, and it
is achieved by \emph{every} horizontally aligned pair (e.g.\ $(3,0)$ and
$(0,0)$, or $(3,1)$ and $(0,1)$). So the margin is $3$. Canonically, the
boundary is $x_1 = 1.5$ with $\mathbf{w}=(w_1,0)$; equality at $x_1=3$
gives $3w_1 + b = 1$ and at $x_1 = 0$ gives $-b = 1$, so $b = -1$ and
$w_1 = 2/3$, hence margin $2/\|\mathbf{w}\| = 2/(2/3) = 3$, matching the
hull distance. Note all four points are support vectors: the closest pair
is non-unique even though the optimal plane is unique.
\end{example}

\begin{problem}[Point versus segment]
\label{ex:pointsegment}
Let $A=\{(2,2)\}$ ($y=+1$) and $B=\{(0,0),(0,2)\}$ ($y=-1$). Sketch
$\operatorname{conv}(B)$, the vertical segment from $(0,0)$ to $(0,2)$.
Project the single positive point $(2,2)$ onto that segment to find the
closest hull point $\mathbf{d}^\star$. Deduce the SVM normal
$\mathbf{w}\propto\mathbf{c}^\star-\mathbf{d}^\star$ and the margin
$\|\mathbf{c}^\star-\mathbf{d}^\star\|$, then solve for the canonical
$(\mathbf{w},b)$ and identify which points are support vectors. (Hint:
the foot of the perpendicular from $(2,2)$ to the segment $x_1=0,\,
0\le x_2\le 2$ is $\mathbf{d}^\star = (0,2)$, so $\mathbf{w}\propto(2,0)$,
margin $2$; $(2,2)$ and $(0,2)$ are support vectors, $(0,0)$ is not.)
\end{problem}

% =====================================================================

\section{What gradient descent actually learns: implicit bias}
\label{sec:implicit}

We have a max-margin classifier as an \emph{optimization problem}
\eqref{eq:svm-primal}. But that is not how neural networks (or even
logistic regression) are trained in practice. In practice we pick a loss
function, write down the average loss over the training set, and run
gradient descent, with \emph{no margin or norm term anywhere in the
objective}. So here is a genuinely surprising fact, and the bridge from
Section~\ref{sec:margins} to the rest of the primer: even though nothing
in the objective mentions the margin, gradient descent on the logistic
loss drifts toward the max-margin solution all by itself. This silent
preference is called \emph{implicit bias} \citep{soudry2018}.

\subsection{The setup and the meaning of ``convergence in direction''}

Fold the labels into the inputs by replacing $\mathbf{x}_n$ with
$y_n\mathbf{x}_n$, so we may assume every label is $+1$ and write the
linear predictor as $\mathbf{w}^\top\mathbf{x}_n$. The data are
\emph{linearly separable} if there is some $\mathbf{w}_*$ with
$\mathbf{w}_*^\top\mathbf{x}_n > 0$ for all $n$ -- a weight vector that
gets every point right \citep{soudry2018}. We minimize the empirical loss
\[
  \mathcal{L}(\mathbf{w}) = \sum_{n=1}^N \ell\!\big(\mathbf{w}^\top\mathbf{x}_n\big),
\]
where $\ell$ is, for instance, the logistic loss
$\ell(u) = \log(1+e^{-u})$ or the exponential loss $\ell(u)=e^{-u}$. Both
are positive, smooth, strictly decreasing, and flatten toward $0$ as
$u\to\infty$. Gradient descent takes steps
\[
  \mathbf{w}(t+1) = \mathbf{w}(t) - \eta\,\nabla\mathcal{L}(\mathbf{w}(t)),
  \qquad
  \nabla\mathcal{L}(\mathbf{w}) = \sum_{n=1}^N \ell'\!\big(\mathbf{w}^\top\mathbf{x}_n\big)\mathbf{x}_n,
\]
with a fixed learning rate $\eta>0$ (a small step-size multiplier).

Here is the catch that makes the question subtle. On separable data, to
push the loss to its infimum $0$ the predictor must send every
$\mathbf{w}^\top\mathbf{x}_n\to+\infty$, which forces
$\|\mathbf{w}(t)\|\to\infty$. There is no finite minimizer to converge to.
But for \emph{classification} only the \emph{direction} of $\mathbf{w}$
matters, since the predicted label is $\operatorname{sign}(\mathbf{w}^\top
\mathbf{x})$ and rescaling $\mathbf{w}$ does not change a sign. So the
right question is about the direction.

\begin{definition}[Convergence in direction]
We say $\mathbf{w}(t)$ \emph{converges in direction} to a unit vector
$\mathbf{u}$ if $\|\mathbf{w}(t)\|\to\infty$ while the normalized iterate
settles, $\mathbf{w}(t)/\|\mathbf{w}(t)\| \to \mathbf{u}$.
\end{definition}

\subsection{The linear separable case, proved}

The centerpiece of this section is the theorem of
\citet{soudry2018}: the direction that gradient descent locks onto is the
$L_2$ max-margin (hard-margin SVM) direction. We state it, then give the
two-part argument it rests on -- a convergence lemma and the
``support-vector gradient-domination'' idea that explains \emph{why} the
SVM direction emerges.

\begin{theorem}[Implicit bias of gradient descent, linear case
\citep{soudry2018}]
\label{thm:soudry}
Let the data be linearly separable and $\ell$ be a smooth, decreasing loss
with an exponential tail (the logistic and exponential losses qualify).
For any small enough learning rate $\eta$ and any starting point,
\[
  \mathbf{w}(t) \;=\; \widehat{\mathbf{w}}\,\log t \;+\; \boldsymbol{\rho}(t),
  \qquad \|\boldsymbol{\rho}(t)\| = O(\log\log t),
\]
where $\widehat{\mathbf{w}}$ is the $L_2$ max-margin vector, the solution
of the hard-margin SVM
$\widehat{\mathbf{w}} = \arg\min \|\mathbf{w}\|^2$ s.t.\
$\mathbf{w}^\top\mathbf{x}_n\ge 1$. Consequently
$\mathbf{w}(t)/\|\mathbf{w}(t)\| \to \widehat{\mathbf{w}}/\|\widehat{\mathbf{w}}\|$.
\end{theorem}

We first record that gradient descent does reach zero loss, with the
norm diverging. This is the ``no finite critical point'' lemma of
\citet{soudry2018}.

\begin{lemma}[The loss goes to zero, the norm diverges \citep{soudry2018}]
\label{lem:diverge}
Under the assumptions of Theorem~\ref{thm:soudry},
$\mathcal{L}(\mathbf{w}(t))\to 0$, $\|\mathbf{w}(t)\|\to\infty$, and
$\mathbf{w}(t)^\top\mathbf{x}_n\to+\infty$ for every $n$.
\end{lemma}
\begin{proof}
Take any separator $\mathbf{w}_*$, so $\mathbf{w}_*^\top\mathbf{x}_n>0$ for
all $n$. Dot the gradient with $\mathbf{w}_*$:
\[
  \mathbf{w}_*^\top \nabla\mathcal{L}(\mathbf{w})
  = \sum_{n=1}^N \ell'\!\big(\mathbf{w}^\top\mathbf{x}_n\big)\,
    \big(\mathbf{w}_*^\top\mathbf{x}_n\big).
\]
Every term is a product of $\ell'(\cdot)<0$ (the loss is strictly
decreasing) and $\mathbf{w}_*^\top\mathbf{x}_n>0$, so every term is
\emph{strictly negative}. Hence at any finite $\mathbf{w}$ the sum is
strictly negative and cannot be $0$, which means \emph{there is no finite
critical point}: $\nabla\mathcal{L}(\mathbf{w})\neq\mathbf{0}$ everywhere.
Gradient descent on a smooth loss with a small enough step size is
guaranteed to drive $\nabla\mathcal{L}(\mathbf{w}(t))\to\mathbf{0}$; since
that can only happen ``at infinity,'' the iterate norm diverges,
$\|\mathbf{w}(t)\|\to\infty$. Because $\ell'(u)\to 0$ only as
$u\to+\infty$, the gradient can vanish only if every margin
$\mathbf{w}(t)^\top\mathbf{x}_n\to+\infty$; this forces
$\ell(\mathbf{w}(t)^\top\mathbf{x}_n)\to 0$ for each $n$, so
$\mathcal{L}(\mathbf{w}(t))\to 0$.
\end{proof}

Now the heart of the matter: why is the limiting \emph{direction} the SVM
direction? The argument is what \citet{soudry2018} call gradient
domination by the support vectors.

\begin{proof}[Why the direction is the max-margin direction (sketch of
Theorem~\ref{thm:soudry})]
Use the exponential loss, $\ell(u)=e^{-u}$, for clarity (the logistic
loss behaves the same way to leading order). Suppose, as
Lemma~\ref{lem:diverge}
permits, that $\mathbf{w}(t)/\|\mathbf{w}(t)\|$ approaches some unit
direction, and write the iterate as $\mathbf{w}(t) = g(t)\,\mathbf{w}_\infty
+ \boldsymbol{\rho}(t)$ with $g(t)\to\infty$ along the limit direction
$\mathbf{w}_\infty$ and $\boldsymbol\rho(t)/g(t)\to 0$. The negative
gradient is

\begin{align*}
  -\nabla\mathcal{L}(\mathbf{w}(t))
  &= \sum_{n=1}^N \exp\!\big(-\mathbf{w}(t)^\top\mathbf{x}_n\big)\,\mathbf{x}_n \\
  &= \sum_{n=1}^N \exp\!\big(-g(t)\,\mathbf{w}_\infty^\top\mathbf{x}_n\big)\,
     \exp\!\big(-\boldsymbol\rho(t)^\top\mathbf{x}_n\big)\,\mathbf{x}_n.
\end{align*}

As $g(t)\to\infty$ the factor $\exp(-g(t)\,\mathbf{w}_\infty^\top
\mathbf{x}_n)$ decays geometrically, and it decays \emph{most slowly} for
the points with the \emph{smallest} value of $\mathbf{w}_\infty^\top
\mathbf{x}_n$ -- precisely the points closest to the boundary, the support
vectors $\arg\min_n \mathbf{w}_\infty^\top\mathbf{x}_n$. All other terms
are exponentially smaller and become negligible. So in the limit the
update direction is dominated by, and is a \emph{nonnegative combination
of}, the support vectors:
\[
  \mathbf{w}_\infty \;\propto\; \widehat{\mathbf{w}}
  \;=\; \sum_{n} \alpha_n \mathbf{x}_n,\qquad
  \alpha_n\ge 0, \ \text{ with } \alpha_n>0 \text{ only for support vectors.}
\]
These are exactly the KKT conditions (Definition~\ref{def:kkt}) of the
hard-margin SVM \eqref{eq:svm-primal}: a nonnegative support-vector
combination, with $\widehat{\mathbf{w}}^\top\mathbf{x}_n=1$ on the
support vectors and $>1$ off them. Hence $\mathbf{w}_\infty$ is
proportional to the SVM solution $\widehat{\mathbf{w}}$. The detailed
proof of \citet{soudry2018} additionally shows $g(t)=\log t$ and bounds
the residual $\boldsymbol\rho(t)$, giving the stated rate.
\end{proof}

Three things to take away, all of which appear later. (1) The convergence
to the SVM direction is \emph{very slow}: the direction error decays like
$O(1/\log t)$, whereas the loss itself decays like $O(1/t)$. So the loss
can be tiny while the margin is still quietly improving -- a reason to keep
training past zero training error, and a warning that loss-based early
stopping can cut off margin (hence robustness) gains \citep{soudry2018}.
(2) The bias is \emph{specific to gradient descent and to exponential-tail
losses}: adaptive methods such as Adam, or polynomial-tail losses, need
not converge to the $L_2$ max-margin direction \citep{soudry2018}. (3) The
same geometric object -- the support vectors, the convex-hull closest
points of Section~\ref{sec:margins} -- is what gradient descent homes in
on.

\subsection{Homogeneous networks: the normalized margin and KKT points}

Linear predictors are a warm-up. Do deep networks, trained the same way,
also drift toward a margin-maximizing solution? \citet{lyuli2020} answer
this for an important class of networks, and we must define two terms
carefully before stating their result.

\begin{definition}[$L$-homogeneous network]
\label{def:homog}
A network $\Phi(\boldsymbol\theta;\mathbf{x})$ with parameters
$\boldsymbol\theta$ is \emph{(positively) $L$-homogeneous} if scaling all
the parameters by a positive constant $c$ scales the output by $c^L$:
\[
  \Phi(c\,\boldsymbol\theta;\mathbf{x}) = c^{L}\,\Phi(\boldsymbol\theta;\mathbf{x})
  \qquad\text{for all } c>0.
\]
Fully-connected and convolutional networks with ReLU (or LeakyReLU)
activations and \emph{no bias terms} are homogeneous, with $L$ equal to
the depth (the number of layers) \citep{lyuli2020}. The reason: ReLU
satisfies $\mathrm{ReLU}(c\,z) = c\,\mathrm{ReLU}(z)$ for $c>0$, and each
linear layer multiplies the output by its weight, so scaling the weights
of all $L$ layers multiplies the final output by $c^L$.
\end{definition}

Just as in the linear case, only the direction of $\boldsymbol\theta$
matters for the predicted class, and by homogeneity the raw margin scales
with $\|\boldsymbol\theta\|^L$. So to compare directions fairly we
normalize.

\begin{definition}[Normalized margin \citep{lyuli2020}]
Let $q_n(\boldsymbol\theta) = y_n\Phi(\boldsymbol\theta;\mathbf{x}_n)$ be
the margin on point $n$, and $q_{\min}(\boldsymbol\theta)
= \min_n q_n(\boldsymbol\theta)$ the smallest margin over the data. The
\emph{normalized margin} divides out the scale,
\[
  \bar\gamma(\boldsymbol\theta)
  = \frac{q_{\min}(\boldsymbol\theta)}{\|\boldsymbol\theta\|^{L}}.
\]
It is invariant to rescaling $\boldsymbol\theta$, just as the geometric
margin was for the linear model.
\end{definition}

\begin{theorem}[Gradient descent maximizes the normalized margin
\citep{lyuli2020}]
\label{thm:lyuli}
Consider an $L$-homogeneous network with an exponential-tail loss
(logistic, exponential, cross-entropy), trained by gradient descent /
gradient flow, and suppose at some time the network reaches $100\%$
training accuracy. Then:
\begin{enumerate}[label=(\alph*),nosep]
\item a smoothed version of the normalized margin $\bar\gamma$ is
  \emph{monotonically non-decreasing} thereafter; and
\item every limit point of the normalized direction
  $\boldsymbol\theta(t)/\|\boldsymbol\theta(t)\|$ lies along a
  \emph{KKT point} of the margin-maximization program
  \[
    \min \ \tfrac12\|\boldsymbol\theta\|^2
    \quad\text{s.t.}\quad q_n(\boldsymbol\theta)\ge 1 \ \ \forall n.
  \]
\end{enumerate}
\end{theorem}

We state Theorem~\ref{thm:lyuli} \emph{without proof}: the argument uses
the Clarke subdifferential (a generalized gradient for the non-smooth ReLU
network), an Euler-type radial/tangential decomposition of the parameter
flow, and a monotonicity argument for the smoothed margin -- machinery
well beyond a first course, and not needed to use the result
\citep{lyuli2020}. \citet{jitelgarsky2019} prove a closely related
statement and, for general (possibly non-separable) data, show gradient
descent follows a unique ray whose direction is the max-margin separator
of the separable part of the data.

\paragraph{An honest caveat: which KKT point?}
The crucial difference from the linear case is the word ``KKT.'' For the
convex SVM, a KKT point \emph{is} the global max-margin solution
(Definition~\ref{def:kkt}). For a non-convex deep network the
margin-maximization program is non-convex, so a KKT point is only a
\emph{first-order stationary point}; it need not be the global
maximum-margin solution, and \citet{lyuli2020} note one can construct
examples where it is not. Theorem~\ref{thm:lyuli} therefore tells us
gradient descent maximizes margin \emph{locally / in a stationarity
sense}, but it does \emph{not} pin down which margin-stationary solution
among possibly many. This honesty matters for the project: ``gradient
descent maximizes the margin'' is exactly true for linear models and only
true ``up to KKT'' for deep nets.

\paragraph{The margin/robustness bridge.}
Why care about margin at all? Because the (normalized) margin
\emph{lower-bounds the distance to the decision boundary}, and that
distance is the robust radius. Concretely \citet{lyuli2020} state that for
a ReLU network the normalized margin on a point, divided by the network's
Lipschitz constant, is a lower bound on the $L_2$ distance from that point
to the boundary (the same Lipschitz--margin certificate that the
robustness section of this primer proves via \citet{tsuzuku2018lipschitz}).
So: plain gradient descent grows the margin, a larger margin means a larger
certified robust radius, and that is the thread connecting ``what training
does'' to ``how robust the model is.''

\begin{example}[2D logistic gradient descent finds the SVM direction]
\label{ex:gd2d}
Take two positive points $\mathbf{x}_1=(1,1)$ and $\mathbf{x}_2=(1,-1)$
(both label $+1$, so no label folding needed) and the exponential loss.
The hard-margin SVM solution is $\widehat{\mathbf{w}}=(1,0)$: it is the
shortest $\mathbf{w}$ with $\mathbf{w}^\top\mathbf{x}_1\ge1$ and
$\mathbf{w}^\top\mathbf{x}_2\ge1$, and by symmetry it points along the
$x_1$-axis. Now run gradient flow from $\mathbf{w}(0)=(0,c)$. The negative
gradient is
\[
  -\nabla\mathcal{L}(\mathbf{w})
  = e^{-(w_1+w_2)}(1,1) + e^{-(w_1-w_2)}(1,-1),
\]
whose second component is $e^{-(w_1+w_2)} - e^{-(w_1-w_2)}$. When $w_2>0$
this is negative (the first exponent is larger, so its term is smaller),
pushing $w_2$ down; when $w_2<0$ it is positive, pushing $w_2$ up. Either
way $w_2\to0$, while the first component is always positive so $w_1\to
+\infty$. Hence $\mathbf{w}(t)/\|\mathbf{w}(t)\|\to(1,0)$: gradient descent
recovers the SVM direction, exactly as Theorem~\ref{thm:soudry} promises.
\end{example}

\begin{problem}[Non-support points do not move the limit]
\label{ex:nonsupport}
Add a third positive point $\mathbf{x}_3=(3,0)$ to
Example~\ref{ex:gd2d}. Check that $\widehat{\mathbf{w}}^\top\mathbf{x}_3
= 3 > 1$, so $\mathbf{x}_3$ is \emph{not} a support vector and its KKT
multiplier is $\alpha_3=0$. Verify the SVM stationarity
$\widehat{\mathbf{w}} = \alpha_1\mathbf{x}_1 + \alpha_2\mathbf{x}_2$, and
argue that although $\mathbf{x}_3$ tugs on the early gradient, the limiting
direction is unchanged and would be the same if $\mathbf{x}_3$ were
deleted. (This is gradient domination by support vectors in action.)
\end{problem}

% =====================================================================

\section{The neural tangent kernel and the lazy-versus-rich dichotomy}
\label{sec:ntk}

The implicit-bias results tell us a network's training drifts toward a
margin-maximizing direction. But \emph{what features} does it use to get
there? Does it learn new features tuned to the task, or does it ride a
fixed set of features baked in at the start? This is the lazy-versus-rich
question, and the tool for making it precise is the \emph{neural tangent
kernel} \citep{jacot2018}. We first build the kernel from a Taylor
expansion, state Jacot's two theorems, give Chizat's criterion for when
training is ``lazy,'' and close with the one definition (feature
averaging) that links all of this back to robustness.

\subsection{Linearization and the neural tangent kernel}

Write a network's scalar output as $f(\boldsymbol\theta, \mathbf{x})$, a
function of parameters $\boldsymbol\theta$ and input $\mathbf{x}$.
Training changes $\boldsymbol\theta$; suppose it stays near its starting
value $\boldsymbol\theta_0$. A first-order Taylor expansion in the
parameters (not the input) gives the \emph{linearization}.

\begin{definition}[Linearization in parameters]
The \emph{linearized model} around $\boldsymbol\theta_0$ is the first-order
Taylor approximation in $\boldsymbol\theta$,
\[
  f(\boldsymbol\theta, \mathbf{x}) \;\approx\;
  f(\boldsymbol\theta_0, \mathbf{x})
  + \nabla_{\boldsymbol\theta} f(\boldsymbol\theta_0, \mathbf{x})
    \cdot (\boldsymbol\theta - \boldsymbol\theta_0).
\]
This is \emph{linear in the parameter shift}
$\boldsymbol\theta-\boldsymbol\theta_0$, with a fixed \emph{feature map}
$\boldsymbol\phi(\mathbf{x}) =
\nabla_{\boldsymbol\theta} f(\boldsymbol\theta_0, \mathbf{x})$ (the
gradient of the output with respect to the parameters, frozen at
$\boldsymbol\theta_0$). It is still nonlinear in $\mathbf{x}$, because
$\boldsymbol\phi(\mathbf{x})$ can be a complicated function of
$\mathbf{x}$.
\end{definition}

\begin{definition}[Neural tangent kernel \citep{jacot2018}]
The \emph{neural tangent kernel} (NTK) is the inner product of two such
feature maps,
\[
  K(\mathbf{x}, \mathbf{x}')
  = \big\langle \nabla_{\boldsymbol\theta} f(\boldsymbol\theta, \mathbf{x}),\,
    \nabla_{\boldsymbol\theta} f(\boldsymbol\theta, \mathbf{x}') \big\rangle
  = \sum_{p} \frac{\partial f}{\partial \theta_p}(\boldsymbol\theta,\mathbf{x})\,
    \frac{\partial f}{\partial \theta_p}(\boldsymbol\theta,\mathbf{x}'),
\]
summing over all parameters $\theta_p$. A \emph{kernel} is just a symmetric
function $K(\mathbf{x},\mathbf{x}')$ that measures similarity between
inputs by an inner product of feature vectors; it lets us do linear
algebra in a feature space without writing the features out. The NTK is the
kernel whose features are the parameter-gradients of the network.
\end{definition}

Under gradient flow the network's \emph{outputs} on the training set move
according to this kernel. If $u_i(t) = f(\boldsymbol\theta(t),\mathbf{x}_i)$
collects the outputs and $\mathbf{y}$ the targets, then
\citep{jacot2018}
\[
  \frac{d\,\mathbf{u}(t)}{dt} = -\,H(t)\,\big(\mathbf{u}(t) - \mathbf{y}\big),
  \qquad H(t)_{ij} = K(\mathbf{x}_i, \mathbf{x}_j; \boldsymbol\theta(t)),
\]
i.e.\ \emph{kernel gradient descent in function space} driven by the NTK
Gram matrix $H$. Two facts make this useful: the cost is \emph{convex in
function space} (even though it is wildly non-convex in
$\boldsymbol\theta$), and the dynamics relax fastest along the top
eigenvectors (principal components) of $H$ -- which is a clean theoretical
motivation for early stopping \citep{jacot2018}.

\subsection{Width, initialization scale, and Jacot's theorems}

Two quantities decide whether the kernel picture is exact: \emph{width}
and \emph{initialization scale}.

\begin{definition}[Width and initialization scale]
The \emph{width} of a hidden layer is its number of neurons; ``infinite
width'' means letting every hidden layer's neuron count grow without
bound. The \emph{initialization scale} is the overall size of the random
parameters at the start, controlled by a factor $\alpha>0$ multiplying the
output (or equivalently the variance of the initial weights). Both are
``knobs'' that, turned up, make the network behave more like its fixed
linearization.
\end{definition}

Jacot's two theorems say that in the infinite-width limit the NTK becomes
deterministic and frozen \citep{jacot2018}.

\begin{theorem}[NTK at initialization \citep{jacot2018}, stated]
\label{thm:jacot1}
As all hidden widths $\to\infty$, the NTK at initialization converges (in
probability) to a deterministic limiting kernel $\Theta_\infty$ that
depends only on the choice of nonlinearity, the depth, and the
initialization variance -- not on the particular random draw of weights.
\end{theorem}

\begin{theorem}[NTK stays constant during training \citep{jacot2018},
stated]
\label{thm:jacot2}
Under mild smoothness, in the same infinite-width limit the NTK stays equal
to $\Theta_\infty$ \emph{throughout} training. Hence training is exactly
kernel gradient descent with one fixed kernel; for a least-squares loss the
network function follows a linear differential equation and the outputs
relax as $\mathbf{u}(t) = \mathbf{y} + e^{-H t}(\mathbf{u}(0)-\mathbf{y})$.
\end{theorem}

We state Theorems~\ref{thm:jacot1}--\ref{thm:jacot2} \emph{without proof}:
they rest on a central-limit-theorem argument over the many neurons and on
controlling how the kernel drifts during training, both of which are
genuinely involved \citep{jacot2018}. The intuition, though, is simple and
worth internalizing: \emph{a wide network has so many neurons that each one
only has to move a tiny amount to fit the data, and if the parameters move
only a little, the first-order Taylor approximation (the linearization)
stays accurate.} One honest subtlety from \citet{jacot2018} (their
Remark~4): each individual neuron's contribution to the kernel does shrink
with width, but the \emph{collective} contribution of the lower layers does
not, which is why even ``lazy'' lower layers still participate.

\subsection{Lazy versus rich, and Chizat's criterion}

We can now name the two regimes. Both terms are now standard, traceable to
\citet{jacot2018} and \citet{chizat2019}.

\begin{definition}[Lazy and rich training]
Training is \emph{lazy} (also called the \emph{kernel} or
\emph{fixed-features} regime) when the NTK barely changes, so the network
behaves like its linearization around $\boldsymbol\theta_0$: it is doing
kernel regression with features \emph{fixed at initialization}, and the
parameters hardly move \citep{chizat2019}. Training is \emph{rich} (the
\emph{feature-learning} regime) when the NTK \emph{evolves} during
training, so the features adapt to the task; this regime lives outside
kernel theory (it is described by mean-field / optimal-transport analyses).
\end{definition}

\citet{chizat2019} pinned down \emph{when} training is lazy with one
dimensionless number. Introduce an explicit scale factor $\alpha$ via the
rescaled objective $F_\alpha(\mathbf{w}) = \tfrac{1}{\alpha^2}
R(\alpha\,h(\mathbf{w}))$. Their criterion (for the square loss $R$) is

\begin{equation}
\label{eq:kappa}
  \kappa \;=\; \frac{\|h(\mathbf{w}_0)-\mathbf{y}^\star\|\;\|D^2 h(\mathbf{w}_0)\|}
                    {\|D h(\mathbf{w}_0)\|^2} \;\ll\; 1
  \quad\Longrightarrow\quad \text{training is lazy},
\end{equation}

where $D h$ and $D^2 h$ are the first and second derivatives of the model
in the parameters. Small $\kappa$ means the model's curvature ($D^2 h$) is
negligible relative to its slope ($Dh$) over the distance training has to
travel, so the linearization holds. The key scaling facts
\citep{chizat2019}:
\begin{itemize}[nosep]
\item \emph{Initialization scale.} Under rescaling by $\alpha$,
  $\kappa_{\alpha h} = \kappa / \alpha$. So \emph{turning up the scale
  $\alpha$ forces laziness}: large $\alpha$ $\Rightarrow$ small $\kappa$
  $\Rightarrow$ lazy. This is the scaling law $\kappa_\alpha = \kappa/\alpha$
  in Chizat et al.\ \citep{chizat2019}.
\item \emph{Width.} For a one-hidden-layer net the gradient norm grows like
  $\sqrt{m}$ (with width $m$) while the curvature stays $O(1)$, so
  $\kappa\sim O(1)/(\sqrt{m})^2 = O(1/m)\to 0$. Width is a second route to
  the same small-$\kappa$ laziness.
\item \emph{Parameters barely move.} In the lazy regime
  $\sup_t\|\mathbf{w}_\alpha(t)-\mathbf{w}_0\| = O(1/\alpha)$, confirming
  ``the parameters hardly change.''
\end{itemize}

\begin{remark}[Chizat's honest counterpoint]
\citet{chizat2019} stress that lazy training is \emph{not} the secret of
deep learning's success. Their experiments show lazy-trained deep CNNs
\emph{underperform} ordinary (rich) training. Laziness gives clean theory
and guaranteed fast fitting, but at the cost of recovering a fixed-kernel
method; the practical wins come from feature learning. We carry this
caveat into the bridge: the lazy/rich distinction is a useful conceptual
axis, not a claim that real networks are lazy.
\end{remark}

\subsection{Worked examples: NTK of simple models}

\begin{example}[NTK of a linear model is constant]
\label{ex:ntklinear}
Let $f(\boldsymbol\theta,\mathbf{x}) = \boldsymbol\theta\cdot\mathbf{x}$
(plain linear model). Then $\nabla_{\boldsymbol\theta} f = \mathbf{x}$ for
\emph{every} $\boldsymbol\theta$, so
\[
  K(\mathbf{x},\mathbf{x}') = \langle \mathbf{x}, \mathbf{x}'\rangle
  = \mathbf{x}\cdot\mathbf{x}',
\]
which does not depend on $\boldsymbol\theta$ or on time at all. A linear
model is ``perfectly lazy'': its NTK is the constant kernel
$\mathbf{x}\cdot\mathbf{x}'$, and training it \emph{is} kernel / linear
regression. (In the scalar case $f=\theta x$ the kernel is just $xx'$.)
\end{example}

\begin{example}[NTK of one hidden unit depends on initialization]
\label{ex:ntkunit}
Let $f(\boldsymbol\theta,x) = b\,\sigma(a x)$ with parameters
$\boldsymbol\theta=(a,b)$ and a smooth activation $\sigma$. Then
$\nabla_{\boldsymbol\theta} f = \big(b\,\sigma'(ax)\,x,\ \sigma(ax)\big)$,
so
\[
  K(x,x') = b^2\,\sigma'(ax)\,\sigma'(ax')\,x x' \;+\; \sigma(ax)\,\sigma(ax').
\]
Unlike the linear case, this \emph{depends on $(a,b)$} -- so at random
initialization the kernel is random, and at finite width it drifts as
$(a,b)$ move during training. What infinite width buys is exactly the
freezing of this dependence into a deterministic $\Theta_\infty$
(Theorem~\ref{thm:jacot1}).
\end{example}

\begin{problem}[Rescaled-model laziness]
\label{ex:rescaled}
Consider a rescaled model $f_\alpha = \alpha\, g$ with the model arranged so
that $g(\mathbf{w}_0,\mathbf{x})=0$ at initialization. Using the criterion
\eqref{eq:kappa}, show $\kappa_{\alpha g} = \tfrac{1}{\alpha}\,\kappa_g$,
and conclude that the linearization error vanishes as $\alpha\to\infty$, so
the model trains lazily. (Extension: verify $\kappa\sim 1/m$ for a
one-hidden-layer net using the $\sqrt{m}$-versus-$O(1)$ scaling of the
gradient and curvature norms.)
\end{problem}

\subsection{Feature averaging and orbit-structured data}

The last definition is the mechanism that ties lazy/rich to the project's
real question: \emph{why a max-margin solution can fail to be robust}.

\begin{definition}[Feature averaging, informally; orbit-structured data]
\label{def:featureavg}
Many natural inputs are built from \emph{shifted (or otherwise
transformed) copies} of a few underlying patterns. Calling the set of all
shifts of a pattern its \emph{orbit} (the group-action term made precise in
the symmetry section of this primer), we say data are
\emph{orbit-structured} when each class is a union of such orbits: a cat is
a cat wherever it sits in the image. \emph{Feature averaging} is the
phenomenon where a network, instead of detecting the discriminative
pattern wherever it appears, learns a weight that effectively sums or
averages the signal across all positions of the orbit. Such an averaged
feature can separate the training data with a large margin, yet a small,
carefully placed perturbation that disturbs the averaging can flip the
prediction -- so the large-margin solution is \emph{not} robust
\citep{li2025feature}.
\end{definition}

We keep Definition~\ref{def:featureavg} at the level of motivation,
because none of the five learning-theory sources read here
(\citealp{soudry2018, lyuli2020, jitelgarsky2019, jacot2018, chizat2019})
proves a direct theorem of the form ``lazy/rich $\Leftrightarrow$
shift-invariant/robust.'' What they \emph{do} give us is the scaffolding:
a lazy network cannot learn a new invariance that is not already in its
initialization kernel (its features are frozen), whereas a rich network
can adapt its features and \emph{learn} shift-invariance rather than
inherit it. Whether the network's margin-maximizing solution is a robust,
invariance-respecting one or a brittle, feature-averaging one therefore
depends on which regime it trained in -- the conceptual hinge the project
investigates.

% =====================================================================

\subsection*{Closing bridge: how the Part-1 thread fits together}

\begin{quote}
\textbf{The one-paragraph story this whole part has been building.}
\emph{Shift-invariance} is, at bottom, a \emph{group-averaging} operation:
averaging a signal over all the shifts in its orbit is the Reynolds
projection onto the invariant subspace (the symmetry section of this
primer makes this an honest orthogonal projection when the shifts act by
orthogonal matrices). The classical \emph{descriptors} of that invariant
content are the \emph{power spectrum} and \emph{bispectrum} (the
shift-invariants section): a shift only rotates the phase of each Fourier
coefficient, leaving $|X|^2$ untouched, and the bispectrum keeps enough
relative phase to tell apart signals the power spectrum confuses.
\emph{Robustness} of a classifier is governed by two numbers we met here
and in the Lipschitz section: the \emph{margin} (Section~\ref{sec:margins})
and the \emph{Lipschitz constant}, with the certified robust radius
behaving like margin over Lipschitz constant. The reason margin is the
right currency is the implicit-bias result
(Section~\ref{sec:implicit}): \emph{plain gradient descent, with no
robustness term in the objective, already drifts toward the max-margin
solution} -- exactly the SVM solution / convex-hull closest points for
linear models, and a KKT point of the margin program for homogeneous
networks. Finally, whether the network \emph{learns} a shift-invariant,
robust representation or merely \emph{inherits} a fixed set of features
(and possibly a brittle, feature-averaging margin) depends on the
\emph{lazy-versus-rich} regime (Section~\ref{sec:ntk}): a lazy,
wide-or-large-scale network is stuck with its initialization kernel and
cannot learn a new invariance, whereas a rich network can. That is the
chain -- \emph{group-averaging $\rightarrow$ spectral descriptors
$\rightarrow$ margin and Lipschitz $\rightarrow$ implicit bias of gradient
descent $\rightarrow$ lazy vs.\ rich} -- on which the rest of the project
is built.
\end{quote}

% =====================================================================
% \citep keys used in this file (to be provided by the host bibliography):
%   vapnik           -- Vapnik 1995/1998, Statistical Learning Theory (max-margin / SVM origin)
%   boser1992        -- Boser, Guyon, Vapnik 1992 (original SVM)
%   cortes1995       -- Cortes, Vapnik 1995 (soft-margin SVM)
%   burges1998       -- Burges 1998 tutorial (canonical margin build-up; asserts distance formula)
%   bennett2000      -- Bennett, Bredensteiner 2000 (duality/geometry, convex-hull closest points)
%   soudry2018       -- Soudry et al. 2018 JMLR (implicit bias, linear separable max-margin)
%   lyuli2020        -- Lyu, Li 2020 ICLR (homogeneous nets, normalized margin -> KKT)
%   jitelgarsky2019  -- Ji, Telgarsky 2019 (non-separable ray / risk-margin)
%   jacot2018        -- Jacot, Gabriel, Hongler 2018 NeurIPS (NTK; Thms 1-2)   [bib key in repo: jacot2018neural]
%   chizat2019       -- Chizat, Oyallon, Bach 2019 NeurIPS (lazy training, kappa criterion)
%   li2025feature    -- feature-averaging non-robustness (used only for the motivation in Def. of feature averaging)  [bib key in repo: li2025feature]
%   tsuzuku2018lipschitz -- Tsuzuku et al. 2018 (Lipschitz-margin certificate; referenced in the bridge)  [bib key in repo]
% Note: jacot2018, li2025feature and tsuzuku2018 already exist in paper/report/references.bib
%       as jacot2018neural, li2025feature, tsuzuku2018lipschitz. If this fragment is
%       compiled against that bib, rename jacot2018 -> jacot2018neural. The SVM/implicit-bias
%       keys (vapnik, boser1992, cortes1995, burges1998, bennett2000, soudry2018, lyuli2020,
%       jitelgarsky2019, chizat2019) are NOT yet in references.bib and must be added.
% =====================================================================
